\chapter {Anhang}

\section{Franck-Hertz-Versuch}

\section{Zeeman-Effekt}

\subsection{Kalibrationsmessungen}

\begin{table}[H]
	\centering
	\begin{tabular}{|c|c|} \hline
		$Ort x / m$         & $B / T$           \\ \hline
		\num{18.5(0.01)e-2} & \num{144.0(1)e-3} \\
		\num{18.6(0.01)e-2} & \num{151.0(1)e-3} \\
		\num{18.7(0.01)e-2} & \num{163.0(1)e-3} \\
		\num{18.8(0.01)e-2} & \num{178.0(1)e-3} \\
		\num{18.9(0.01)e-2} & \num{198.0(1)e-3} \\
		19.0e-2, 198.0e-3
		19.1e-2, 288.0e-3
		19.2e-2, 322.0e-3
		19.3e-2, 368.0e-3
		19.4e-2, 420.0e-3
		19.5e-2, 455.0e-3
		19.6e-2, 466.0e-3
		19.7e-2, 466.0e-3
		19.8e-2, 458.0e-3
		19.9e-2, 453.0e-3
		20.0e-2, 458.0e-3
		20.1e-2, 470.0e-3
		20.2e-2, 479.0e-3
		20.3e-2, 477.0e-3
		20.4e-2, 454.0e-3
		20.5e-2, 414.0e-3
		20.6e-2, 371.0e-3
		20.7e-2, 324e-3
		20.8e-2, 284e-3
		20.9e-2, 259e-3
		21.0e-2, 225e-3
		21.1e-2, 202e-3
		21.2e-2, 184e-3
		21.3e-2, 170e-3
		21.4e-2, 155e-3
		21.5e-2, 143e-3
		\hline
	\end{tabular}
	\caption{Messwerte der Ortskalibration der Magnetfeldstärke bei $I=\SI{5.47(0.01)}{\ampere}$.}
	\label{tab:ortskalibration}
\end{table}
